% RESTful APIs & SignalR section

% 1 – REST Constraints
\QuestionSlide[\CategoryBadge[RestColor!20]{REST \& SignalR}]<1>{What are the six REST architectural constraints?}

\begin{frame}[fragile]
  \frametitle{\AnswerTitle[\CategoryBadge[RestColor!20]{REST \& SignalR}]{The six REST constraints?}}

  {\footnotesize
  \begin{enumerate}
    \item \textbf{Client-Server} – UI and data store are separated.
    \item \textbf{Stateless} – each request carries all context; no server-side session.
    \item \textbf{Cacheable} – responses declare cacheability via headers.
    \item \textbf{Uniform Interface} – resource-based URIs, standard HTTP verbs, self-descriptive messages.
    \item \textbf{Layered System} – intermediaries (proxies, gateways) are transparent.
    \item \textbf{Code on Demand} (optional) – server can send executable code to client.
  \end{enumerate}
  }
\end{frame}

% 2 – HTTP Verbs & Status Codes
\QuestionSlide[\CategoryBadge[RestColor!20]{REST \& SignalR}]<1>{What HTTP verbs and status codes are standard in REST APIs?}

\begin{frame}[fragile]
  \frametitle{\AnswerTitle[\CategoryBadge[RestColor!20]{REST \& SignalR}]{HTTP verbs and status codes?}}

  {\footnotesize
  \begin{tabular}{lll}
    \textbf{Verb} & \textbf{Action} & \textbf{Success Code} \\
    GET    & Read          & 200 OK \\
    POST   & Create        & 201 Created + Location header \\
    PUT    & Full replace  & 200 OK / 204 No Content \\
    PATCH  & Partial update & 200 / 204 \\
    DELETE & Remove        & 204 No Content \\
  \end{tabular}
  \vspace{0.5em}
  Key error codes: 400 Bad Request, 401 Unauthorized, 403 Forbidden, 404 Not Found, 409 Conflict, 422 Unprocessable Entity, 429 Too Many Requests.
  }
\end{frame}

% 3 – API Versioning
\QuestionSlide[\CategoryBadge[RestColor!20]{REST \& SignalR}]<2>{What are the main strategies for REST API versioning?}

\begin{frame}[fragile]
  \frametitle{\AnswerTitle[\CategoryBadge[RestColor!20]{REST \& SignalR}]{REST API versioning strategies?}}

  {\footnotesize
  \begin{itemize}
    \item \textbf{URI path}: \texttt{/api/v1/products} — most visible, easy to cache.
    \item \textbf{Query string}: \texttt{/api/products?api-version=1.0} — no URL changes.
    \item \textbf{Header}: \texttt{X-Api-Version: 1.0} — clean URLs, less discoverable.
    \item \textbf{Media type (Accept)}: \texttt{application/vnd.myapi.v1+json} — purist REST, complex clients.
  \end{itemize}
  Use \textbf{Asp.Versioning.Mvc} NuGet in ASP.NET Core for any strategy.
  }

  \begin{minted}[fontsize=\scriptsize]{csharp}
builder.Services.AddApiVersioning(opt =>
{
    opt.DefaultApiVersion = new ApiVersion(1, 0);
    opt.AssumeDefaultVersionWhenUnspecified = true;
    opt.ReportApiVersions = true;
});
  \end{minted}
\end{frame}

% 4 – OpenAPI / Swagger
\QuestionSlide[\CategoryBadge[RestColor!20]{REST \& SignalR}]<1>{What is OpenAPI / Swagger in ASP.NET Core?}

\begin{frame}[fragile]
  \frametitle{\AnswerTitle[\CategoryBadge[RestColor!20]{REST \& SignalR}]{What is OpenAPI/Swagger?}}

  {\footnotesize
  \textbf{OpenAPI} is a machine-readable JSON/YAML specification for REST APIs. \textbf{Swagger UI} renders it as interactive docs. In ASP.NET Core, \texttt{Microsoft.AspNetCore.OpenApi} (built-in .NET 9) or \textbf{Swashbuckle} generates the spec from controller metadata and XML comments. Clients can auto-generate SDKs with \texttt{NSwag} or \texttt{kiota}.
  }

  \begin{minted}[fontsize=\scriptsize]{csharp}
// .NET 9 built-in
builder.Services.AddOpenApi();
app.MapOpenApi();  // serves /openapi/v1.json

// With Swashbuckle (older)
builder.Services.AddSwaggerGen(c =>
    c.SwaggerDoc("v1", new OpenApiInfo { Title = "MyApi", Version = "v1" }));
app.UseSwagger();
app.UseSwaggerUI();
  \end{minted}
\end{frame}

% 5 – Problem Details
\QuestionSlide[\CategoryBadge[RestColor!20]{REST \& SignalR}]<2>{What is Problem Details (RFC 7807) in ASP.NET Core?}

\begin{frame}[fragile]
  \frametitle{\AnswerTitle[\CategoryBadge[RestColor!20]{REST \& SignalR}]{What is Problem Details?}}

  {\footnotesize
  \textbf{Problem Details} (RFC 7807) is a standard JSON format for HTTP error responses: \texttt{type}, \texttt{title}, \texttt{status}, \texttt{detail}, and \texttt{instance}. ASP.NET Core \texttt{[ApiController]} returns Problem Details automatically for 4xx errors. \texttt{IProblemDetailsService} lets you customise and extend the format.
  }

  \begin{minted}[fontsize=\scriptsize]{json}
{
  "type": "https://tools.ietf.org/html/rfc7231#section-6.5.1",
  "title": "One or more validation errors occurred.",
  "status": 400,
  "traceId": "00-abc123-01",
  "errors": {
    "CustomerName": ["The CustomerName field is required."]
  }
}
  \end{minted}
\end{frame}

% 6 – What is SignalR?
\QuestionSlide[\CategoryBadge[RestColor!20]{REST \& SignalR}]<1>{What is SignalR?}

\begin{frame}[fragile]
  \frametitle{\AnswerTitle[\CategoryBadge[RestColor!20]{REST \& SignalR}]{What is SignalR?}}

  {\footnotesize
  \textbf{SignalR} is ASP.NET Core's real-time communication library. It abstracts transport selection: \textbf{WebSockets} (preferred), \textbf{Server-Sent Events}, and \textbf{Long Polling} fallbacks. The server \textbf{pushes} messages to clients rather than clients polling. Used in live dashboards, chat, notifications, multiplayer games, and Blazor Server.
  }
\end{frame}

% 7 – SignalR Hubs
\QuestionSlide[\CategoryBadge[RestColor!20]{REST \& SignalR}]<2>{What are SignalR Hubs?}

\begin{frame}[fragile]
  \frametitle{\AnswerTitle[\CategoryBadge[RestColor!20]{REST \& SignalR}]{What are SignalR Hubs?}}

  {\footnotesize
  A \textbf{Hub} is a server-side class that acts as the central communication endpoint. Clients call Hub methods; the Hub calls client methods via \texttt{Clients.All}, \texttt{Clients.Caller}, or \texttt{Clients.Group(name)}. Hubs are transient — no state between calls; use DI for persistent services.
  }

  \begin{minted}[fontsize=\scriptsize]{csharp}
public class ChatHub : Hub
{
    public async Task SendMessage(string user, string message)
        => await Clients.All.SendAsync("ReceiveMessage", user, message);

    public async Task JoinGroup(string group)
        => await Groups.AddToGroupAsync(Context.ConnectionId, group);
}

// Registration
app.MapHub<ChatHub>("/chat");
  \end{minted}
\end{frame}

% 8 – WebSockets vs SSE vs Long Polling
\QuestionSlide[\CategoryBadge[RestColor!20]{REST \& SignalR}]<2>{When would you choose WebSockets vs SSE vs Long Polling?}

\begin{frame}[fragile]
  \frametitle{\AnswerTitle[\CategoryBadge[RestColor!20]{REST \& SignalR}]{WebSockets vs SSE vs Long Polling?}}

  {\footnotesize
  \begin{tabular}{p{2.3cm}p{2.3cm}p{2.3cm}}
    \textbf{WebSockets} & \textbf{SSE} & \textbf{Long Polling} \\
    Full-duplex & Server $\to$ client only & Client polls with hanging request \\
    Low overhead & HTTP/2 multiplexed & Highest overhead \\
    Chat, games, Blazor & Live feeds, notifications & Fallback / firewalls \\
  \end{tabular}
  \vspace{0.5em}
  SignalR negotiates from WebSockets $\to$ SSE $\to$ Long Polling automatically per client capability.
  }
\end{frame}
